\section{Analysis}
\label{sec:analysis}

\paragraph{BM25 top-k tradeoff.}
The current results show that sequential top-100 is not consistently the strongest setting on the 50-query subsets.
Top-80 gives the highest NDCG@10 on SciFact, SCIDOCS, and FiQA-2018, while top-40 is highest on ArguAna and top-60 is slightly highest on NFCorpus.
This suggests that increasing the BM25 candidate pool to 100 does not automatically improve the early ListT5 reranking result under this small-query setup.
For runtime and forward-call count, top-60 is currently the cheapest setting on all five datasets, while top-80 and top-100 are slower and require more model calls.

\paragraph{Grouping effect.}
Score-balanced grouping performs below sequential top-100 on NFCorpus, SciFact, ArguAna, and SCIDOCS, while producing a nearly tied result on FiQA-2018.
Overall, these preliminary results do not show a clear effectiveness advantage for score-balanced grouping over sequential grouping.


\paragraph{Dataset variation.}
The best BM25 top-k setting differs by dataset, so we avoid relying only on an average score.
For example, ArguAna favors a much smaller top-40 candidate set in the current table, while SCIDOCS and FiQA-2018 favor top-80.
This variation is consistent with the paper's framing as a preliminary diagnostic study rather than a final benchmark result.

\paragraph{Caching effect.}
The strategy-aware cache substantially reduces score-balanced runtime across all five datasets.
Aligned caching is roughly two times faster than using the original sequential-cache behavior.
The forward-call counts show the same trend: aligned caching reduces score-balanced model calls by about two times compared with the original sequential-cache behavior.
However, score-balanced top-100 remains slower than sequential top-100 and much slower than the fastest sequential top-k setting in the current runs.
